% created by Uwe Schadewald
% modified by Mathias Kuntze and Ahmet Uysal
% Add, handout to documentclass arguments for condensed pdf
\documentclass[presentation, 8pt, mathserif, t]{beamer} % , aspectratio=169
\usepackage[english]{babel}
\usepackage{pgf,graphicx}
\usepackage{amsmath, amssymb}
\usepackage[utf8]{inputenc}
\usepackage{lmodern}
\usepackage{palatino}
\usepackage{multimedia}
\usepackage{pgfpages} 
\usepackage{tikz}
\usepackage{datetime}
\pdfoptionpdfminorversion=5

\usepackage{caption}
\usepackage{subcaption}
% if else
\usepackage{ifthen}
% extend table options
\usepackage{tabularx} 
\usepackage{booktabs}
\usepackage{multicol}
\usepackage{multirow}
\usepackage{eso-pic}  % package to set background image
\usepackage[calc]{picture}

% Packages and stuff for ToDo list like itempoints
\usepackage{pifont}
\newcommand{\cmark}{\ding{51}}%
\newcommand{\xmark}{\ding{55}}%
\newcommand{\open}{$\square$}
\newcommand{\done}{\rlap{$\square$}{\raisebox{1pt}{\large\hspace{1.5pt}\cmark}}\hspace{-2.5pt}}
\newcommand{\wontfix}{\rlap{$\square$}{\raisebox{1pt}{\large\hspace{1.5pt}\xmark}}}
\newcommand{\notsure}{\rlap{$\square$}{\raisebox{0.8pt}{\large\hspace{1.5pt}\textbf{?}}}}



% side bar and footer
\setbeamertemplate{headline}{	
	\leavevmode
	\vspace{-4em}	
	\hbox{		
		\begin{beamercolorbox}[wd=0.85\paperwidth,ht=10ex,dp=8ex,center]{}%			
			% navigation with subsections as dots
			\hspace{3.5em}\insertnavigation{0.7\paperwidth}{\hskip0pt plus1fill} % add navigation in footer						
			% navigation with sections, no subsections
			% \insertsectionnavigationhorizontal{0.6\paperwidth}{\hskip0pt plus1fill}{} \\ % add navigation in footer}
			
		\end{beamercolorbox} 				
	}
	\vskip0pt
}


\setbeamertemplate{footline}{	
	\leavevmode
	\vspace{-3em}
	\hbox{
		\begin{beamercolorbox}[wd=.33\paperwidth,ht=2.25ex,dp=1ex,left]{author in head/foot}%
			\hspace{5em}
			\insertshortauthor
		\end{beamercolorbox}
		\begin{beamercolorbox}[wd=.33\paperwidth,ht=2.25ex,dp=1ex,center]{title in head/foot}%
			\insertshorttitle \ - \insertshortsubtitle
		\end{beamercolorbox}	
		\begin{beamercolorbox}[wd=0.30\paperwidth,ht=10ex,dp=8ex,right]{pagenumber in head/foot}			 	
			\insertframenumber % add page numbers
		\end{beamercolorbox}
	}			
	\vskip0pt
}



\setbeamertemplate{frametitle}{
	\ifthenelse{\equal{\insertframesubtitle}{}}{
		\vspace{0.6cm}
		\huge{\insertframetitle}
	}{
		\vspace{0.6cm}
		\small{\insertframetitle}\\
		\vspace{0.3cm}
		\huge{\insertframesubtitle}
    }		
}

	
% enumerate sections
\setbeamertemplate{section in head/foot}{\hfill\insertsectionheadnumber.~\insertsectionhead}
%\setbeamertemplate{section in head/foot shaded}{\color{structure!50}\hfill\insertsectionheadnumber.~\insertsectionhead}
\setbeamertemplate{section in toc}{\inserttocsectionnumber.~\inserttocsection}

%enumerate subsections
\setbeamertemplate{subsection in head/foot}{\hfill\insertsubsectionheadnumber.~\insertsubsectionhead}
\setbeamertemplate{subsection in head/foot shaded}{\color{structure!50}\hfill\insertsubsectionheadnumber.~\insertsubsectionhead}
%\setbeamertemplate{subsection in toc}[subsections numbered]
\setbeamertemplate{subsection in toc}{\vskip0.5em\leftskip=2em\inserttocsubsection\par}

%--------------------------Common------------------------------------------------------
\setbeamercovered{transparent} % make the beamer theme invisible
\usefonttheme{structurebold}
\beamertemplatenavigationsymbolsempty % set navigations helper function to off
\setbeamertemplate{bibliography item}[text]
\setbeamertemplate{note page}[plain]

%\setlist[itemize,1]{label={$\bullet$}} % \item are using bullets
\setbeamertemplate{itemize items}[circle]
	
	

	
% create a new command to show it on two screens
% I'm using dspdfviewer.
\newcommand{\setDualView} {
	\setbeameroption{show notes on second screen=right}
}

%\AtBeginSection[]{\subsection{}}
\newcommand{\addcite}[1]{%
	\AddToShipoutPictureFG*{%
		\AtPageLowerLeft{%
			\put(0.90\paperwidth,5em){											
				\tiny{
					\cite{#1} 
				}			
			}
		}
	}	
}

% insert a frame with references -> use bibtex
\newcommand{\insertReferenceFrame}[3]{%
	\section{#1}
	\begin{frame}[allowframebreaks]
		\frametitle{#1}
		\bibliographystyle{#2}
		\bibliography{#3}
	\end{frame}	
}

\AtBeginSection[]{\subsection{}}
	





\usepackage{../KU-Beamer-Template/style/koc} 
\usepackage{minted}
\usepackage{upquote}
\usepackage{graphicx}

% pdflatex --shell-escape lecture4.tex & pdflatex --shell-escape lecture4.tex

\title{KOLT Python}
% https://www.geeksforgeeks.org/difference-method-function-python/
\subtitle{Lists \& For Loops} 
\newdate{date}{14}{10}{2019}
\date{\displaydate{date}}
\author{Ahmet Uysal}

\titlegraphic{\includegraphics[scale=0.2]{../KU-Beamer-Template/style/images/logo_kolt.pdf}}

\setbeamercovered{invisible} % transparent
\makeatletter
\let\@@magyar@captionfix\relax
\makeatother
\begin{document}
    \maketitle

    \frame{\frametitle{Agenda}\tableofcontents}

    \section{Recap}

        \begin{frame}{Strings}
            \LARGE
            \textbf{\texttt{my\_string = \textquotesingle abcde\textquotesingle}}
            \begin{center}
                \huge
                \begin{table}[]
                    \begingroup
                    \setlength{\tabcolsep}{1pt} % Default value: 6pt
                    \begin{tabular}{rcccccl}
                    & {\color[HTML]{c4122e} 0}  & {\color[HTML]{c4122e} 1}  & {\color[HTML]{c4122e} 2}  & {\color[HTML]{c4122e} 3}\  & {\color[HTML]{c4122e} 4} &  \\
                    \textquotesingle & a  & b  & c  & d  & e & \textquotesingle \\
                    & {\color[HTML]{A0A0A0} -5} & {\color[HTML]{A0A0A0} -4} & {\color[HTML]{A0A0A0} -3} & {\color[HTML]{A0A0A0} -2} & {\color[HTML]{A0A0A0} -1} & 
                    \end{tabular}
                    \endgroup
                    \end{table}
            \end{center}

            \textbf{\texttt{print(my\_string[2])}} $\Rightarrow$ prints c \\
            \textbf{\texttt{print(my\_string[-2])}} $\Rightarrow$ prints d
        \end{frame}

        \begin{frame}{Indexing \& Slicing}
            \Large
            Access specific characters using \textbf{indexing}, i.e, \texttt{\textbf{[index]}}\\
            \textbf{Slice} strings by using \texttt{\textbf{[start:stop:step]}}
            \inputminted[frame=single,framesep=2pt]{python3}{../Lecture3/code-examples/string_index.py}
            \inputminted[frame=single,framesep=2pt]{python3}{../Lecture3/code-examples/string_index2.py}
        \end{frame}

        \begin{frame}{String Operations}
            \large
            \inputminted[frame=single,framesep=2pt, lastline=4]{python3}{../Lecture3/code-examples/string_operations.py}
            \inputminted[frame=single,framesep=2pt, firstline=6]{python3}{../Lecture3/code-examples/string_operations.py}
            \LARGE
            \texttt{\textbf{str1 + str2}} $\Rightarrow$ \textbf{Concatenate} \texttt{str1} and \texttt{str2}\\
            \texttt{\textbf{str1 *}} $n$ $\Rightarrow$  Repeate \texttt{str1} $n$ times.
        \end{frame}

        \begin{frame}{While Loops}
            Repeat some \texttt{<expression>}s \underline{as long as} a \texttt{<condition>} is \texttt{True}.
            \begin{columns}
                \begin{column}{0.45\textwidth}
                    \inputminted[frame=single,framesep=2pt]{python3}{../Lecture3/code-examples/while1.py}
                \end{column}
                \begin{column}{0.45\textwidth}
                    \inputminted[frame=single,framesep=2pt]{python3}{../Lecture3/code-examples/while2.py}
                \end{column} 
            \end{columns}
            \LARGE
            \inputminted[frame=single,framesep=2pt]{python3}{../Lecture3/code-examples/while3.py}
            \texttt{<condition>} is only checked \textbf{\underline{before}} each execution.
        \end{frame}

        \begin{frame}{Lists}
            \begin{center}
                \includegraphics[width=0.6\textwidth]{../Lecture1/images/box_many.jpg}                
            \end{center}
            \LARGE
            Imagine variables, but with limitless capacity$\dots$\\
            \textbf{\texttt{sunnyside = [\textquotesingle Mr. Potato Head\textquotesingle, \textquotesingle Hamm\textquotesingle,
            \textquotesingle Buzz Lightyear\textquotesingle, \textquotesingle Slinky Dog\textquotesingle]}}
        \end{frame}

        \begin{frame}{Lists}
            \huge
            \inputminted[frame=single,framesep=2pt]{python3}{../Lecture3/code-examples/intro_lists.py}
            \inputminted[frame=single,framesep=2pt]{python3}{../Lecture3/code-examples/mixed_list.py}
        \end{frame}

        \begin{frame}{Accessing Elements}
            \LARGE
            \texttt{\textbf{values = [1, \textquotesingle hello\textquotesingle, None, [3], True]}}\\
            \begin{center}
                \huge
                \begin{table}[]
                    \begingroup
                    \setlength{\tabcolsep}{2pt} % Default value: 6pt
                    \begin{tabular}{rcccccl}
                     & {\color[HTML]{c4122e} 0}  & {\color[HTML]{c4122e} 1}  & {\color[HTML]{c4122e} 2}  & {\color[HTML]{c4122e} 3}\  & {\color[HTML]{c4122e} 4} &  \\
                     $[$ & 1,  & \textquotesingle hello\textquotesingle,  & \texttt{None},  & \texttt{[3]},  & \texttt{True} & $]$ \\
                    & {\color[HTML]{A0A0A0} -5} & {\color[HTML]{A0A0A0} -4} & {\color[HTML]{A0A0A0} -3} & {\color[HTML]{A0A0A0} -2} & {\color[HTML]{A0A0A0} -1} & 
                    \end{tabular}
                    \endgroup
                    \end{table}
            \end{center}
            Use \textbf{indexing} to access and \textbf{update} elements inside list.\\
            \textbf{\texttt{print(values[2])}}\\
            \textbf{\texttt{values[2] = \textquotesingle new value\textquotesingle}} 

        \end{frame}

        \begin{frame}{Adding New Elements}
            \LARGE
            Append elements at the end of a list by \textbf{append()}
            \bigskip
            \normalsize
            \inputminted[frame=single,framesep=2pt]{python3}{../Lecture3/code-examples/append_list.py}
            \inputminted[frame=single,framesep=2pt]{python3}{../Lecture3/code-examples/append_list2.py}
        \end{frame}
        
        \begin{frame}{Inspecting List Elements}
            \LARGE
            Slice lists by using \texttt{\textbf{[start:stop:step]}}
            \normalsize
            \begin{columns}
                \begin{column}{0.35\textwidth}
                \inputminted[frame=single,framesep=2pt]{python3}{../Lecture3/code-examples/slicing.py} 
                \end{column}
                \begin{column}{0.65\textwidth}
                 \inputminted[frame=single,framesep=2pt]{python3}{../Lecture3/code-examples/slicing2.py}
                \end{column} 
             \end{columns}
        \end{frame}

        \begin{frame}{Inspecting List Elements}
            \LARGE
            \inputminted[frame=single,framesep=2pt]{python3}{../Lecture3/code-examples/slicing3.py}
            \inputminted[frame=single,framesep=2pt]{python3}{../Lecture3/code-examples/slicing4.py}
        \end{frame}

    \section{Lists(Cont.)}
        \begin{frame}{Removing An Element}
            \pause
            \LARGE
            Remove elements in a list by \textbf{remove()}
            \pause              
            \bigskip
            \normalsize
            \inputminted[frame=single,framesep=2pt]{python3}{../Lecture3/code-examples/remove.py}
            \pause
            \LARGE
            How to avoid \texttt{ValueError}? \pause (Hint: \textbf{Branching})
        \end{frame}

        \begin{frame}{in Operator}
            \pause
            \LARGE
            Search an operand in the specified sequence by using \textbf{in}
            \pause
            \bigskip
            \inputminted[frame=single,framesep=2pt]{python3}{../Lecture3/code-examples/in_operator.py}
            \pause
            \begin{itemize}
                \item Works with both lists and strings
                \pause
                \item Works with ranges
            \end{itemize}
        \end{frame}

        \begin{frame}{len() Function}
            \pause
            \LARGE
            \texttt{len()} is an operator to determine the size of lists, strings, etc.
            \pause
            \bigskip
            \inputminted[frame=single,framesep=2pt]{python3}{../Lecture3/code-examples/length.py}
        \end{frame}

        \begin{frame}{List Slicing}
            \Large
            Access collection of elements with \textbf{\texttt{[start:stop:step]}}\\
            Gives a list, even when number of elements is not bigger than 1.
            \inputminted[frame=single,framesep=2pt,firstline=10]{python3}{../Review1/code-examples/slicing.py}                                
            \LARGE
            Slices with \texttt{step} = 1 are called \textbf{Basic Slice}.\\
            Slices with \texttt{step != 1} are called \textbf{Extended Slice}.
        \end{frame}

        \begin{frame}{List Mutation}
            \Large
            \textbf{\texttt{list.append(x)}}: Append x to end of the sequence\\
            \textbf{\texttt{list.insert(i, x)}}: Insert x to index i\\
            \textbf{\texttt{list.pop(i=-1)}}: Remove and return element at index i\\
            \textbf{\texttt{list.remove(x)}}: Remove first occurrence of x\\
            \textbf{\texttt{list.extend(iterable)}}: Add all elements in iterable to end of list\\
            \textbf{\texttt{list[i] = new\_value}}: Update value of index i with new value\\
            \textbf{\texttt{list[basic\_slice] = iterable}}: Change elements in basic slice with elements in iterable, sizes can be different: \texttt{numbers[:] = []}\\
            \textbf{\texttt{list[extended\_slice] = iterable}}: Change elements in extended slice with elements in iterable 1-1, sizes must be equal.\\
        \end{frame}

        \begin{frame}{Some Other List Operations}
            \Large
            \textbf{\texttt{in}} operator: Check whether an element is in list.\\
            \texttt{3 in numbers} $\Rightarrow$ \texttt{True}\\
            \textbf{\texttt{len(list)}}: Returns the length of list(and other collections).\\
            \textbf{\texttt{list.index(value, start=0, stop=len(list))}}:\\
            Return first index of value.\\
            \textbf{\texttt{list.count(value)}}: Count number of occurrences of value.\\
            \textbf{\texttt{list.reverse()}}: Reverse the list (in-place)\\
            \textbf{\texttt{list.sort()}}: Sort list elements (in-place)\\
            \\ 
            For more, type \texttt{help(list)} in your interactive interpreter.
        \end{frame}

        \section{For Loops}
      
        \begin{frame}{range() Function}
            \pause
            \LARGE
            \texttt{range(\textbf{start, stop, step})} is a function to create ranges
            \bigskip
            \inputminted[frame=single,framesep=2pt]{python3}{../Lecture3/code-examples/range.py}
        \end{frame}
      
        \begin{frame}{For Loops}
            \pause
            \begin{columns}
                \begin{column}{0.5\textwidth}
                    \inputminted[frame=single,framesep=2pt]{python3}{../Lecture3/code-examples/for1.py}
                \end{column}
               \pause 
                \begin{column}{0.5\textwidth}
                    \inputminted[frame=single,framesep=2pt]{python3}{../Lecture3/code-examples/for2.py}
                    \pause
                    \inputminted[frame=single,framesep=2pt]{python3}{../Lecture3/code-examples/for3.py}
                    \pause
                    \inputminted[frame=single,framesep=2pt]{python3}{../Lecture3/code-examples/for4.py}
                \end{column} 
            \end{columns}
        \end{frame}
        
        \begin{frame}{Example: Mail Sender}
            \pause
            \LARGE
            Fill out the attendance form: \href{https://tiny.cc/kolt-python}{\underline{\textit{tiny.cc/kolt-python}}} 
            \begin{center}
                \includegraphics[height=0.65\textheight]{images/unicorn.jpg}                
            \end{center}
        
        \end{frame}

        \begin{frame}{Break, Continue \& Pass}
            \begin{columns}
                \begin{column}{0.5\textwidth}
                    \textbf{\texttt{break}} immidiately terminates the closest loop
                    \bigskip  
                    \inputminted[frame=single,framesep=2pt]{python3}{../Lecture3/code-examples/break1.py}
                    \pause
                    \inputminted[frame=single,framesep=2pt]{python3}{../Lecture3/code-examples/break2.py}
                \end{column}
               \pause 
                \begin{column}{0.5\textwidth}
                    \textbf{\texttt{continue}} skips to the next iteration of the loop
                    \bigskip  
                    \inputminted[frame=single,framesep=2pt]{python3}{../Lecture3/code-examples/continue1.py}
                    \pause
                    \inputminted[frame=single,framesep=2pt]{python3}{../Lecture3/code-examples/continue2.py}
                \end{column} 
            \end{columns}
        \end{frame}
        
        \begin{frame}{Break, Continue \& Pass}
            \pause
            \LARGE
            \textbf{\texttt{pass}} does not have an effect
            \bigskip  
            \inputminted[frame=single,framesep=2pt]{python3}{../Lecture3/code-examples/pass.py}
            \pause
            \begin{itemize}
                \item Loops, conditional statements, functions etc. cannot be empty 
            \end{itemize}
        \end{frame}
\end{document}